\documentclass[UTF8,12pt]{ctexart}
\usepackage[paperwidth=240mm,paperheight=200mm,margin=14mm,top=8mm]{geometry}
\usepackage{amsmath,amssymb,mathtools}
\usepackage{xcolor}
\pagestyle{empty}
\setlength{\parindent}{0pt}
\setlength{\parskip}{0pt}
\newcommand{\qn}[1]{\textbf{\large #1}\ }
\xeCJKDeclareCharClass{CJK}{"2460 -> "24FF}
\begin{document}
\qn{1991.1}
设 $y=\ln(1+3^{-x})$，则 $\mathrm{d}y=$______.

\vfill
\newpage
\qn{1991.2}
曲线 $y=\mathrm{e}^{-x^2}$ 的凸区间是______.

\vfill
\newpage
\qn{1991.3}
$\int_1^{+\infty}\dfrac{\ln x}{x^2}\,\mathrm{d}x=$______.

\vfill
\newpage
\qn{1991.4}
质点以速度 $t\sin(t^2)$ 米/秒作直线运动，则从时刻 $t_1=\sqrt{\dfrac{\pi}{2}}$ 秒到 $t_2=\sqrt{\pi}$ 秒内质点所经过的路程等于______米.

\vfill
\newpage
\qn{1991.5}
$\lim\limits_{x\to0^+}\dfrac{1-\mathrm{e}^{\frac{1}{x}}}{x+\mathrm{e}^{\frac{1}{x}}}=$______.

\vfill
\newpage
\qn{1991.6}
若曲线 $y=x^2+ax+b$ 和 $2y=-1+xy^3$ 在点 $(1,-1)$ 处相切，其中 $a,b$ 是常数，则（　　）
\\
\noindent (A) $a=0,\,b=-2$.　　(B) $a=1,\,b=-3$.
\\
\noindent (C) $a=-3,\,b=1$.　　(D) $a=-1,\,b=-1$.

\vfill
\newpage
\qn{1991.7}
设函数 $f(x)=\begin{cases} x^2, & 0\leqslant x\leqslant 1, \\ 2-x, & 1<x\leqslant 2, \end{cases}$ 记 $F(x)=\int_0^x f(t)\,\mathrm{d}t$，$0\leqslant x\leqslant 2$，则（　　）
\\
\noindent (A) $F(x)=\begin{cases} \dfrac{x^3}{3}, & 0\leqslant x\leqslant 1, \\[6pt] \dfrac{1}{3}+2x-\dfrac{x^2}{2}, & 1<x\leqslant 2. \end{cases}$
\\
\noindent (B) $F(x)=\begin{cases} \dfrac{x^3}{3}, & 0\leqslant x\leqslant 1, \\[6pt] -\dfrac{7}{6}+2x-\dfrac{x^2}{2}, & 1<x\leqslant 2. \end{cases}$
\\
\noindent (C) $F(x)=\begin{cases} \dfrac{x^3}{3}, & 0\leqslant x\leqslant 1, \\[6pt] \dfrac{x^3}{3}+2x-\dfrac{x^2}{2}, & 1<x\leqslant 2. \end{cases}$
\\
\noindent (D) $F(x)=\begin{cases} \dfrac{x^3}{3}, & 0\leqslant x\leqslant 1, \\[6pt] 2x-\dfrac{x^2}{2}, & 1<x\leqslant 2. \end{cases}$

\vfill
\newpage
\qn{1991.8}
设函数 $f(x)$ 在 $(-\infty,+\infty)$ 内有定义，$x_0\neq 0$ 是函数 $f(x)$ 的极大值点，则（　　）
\\
\noindent (A) $x_0$ 必是 $f(x)$ 的驻点.　　(B) $-x_0$ 必是 $-f(-x)$ 的极小值点.
\\
\noindent (C) $-x_0$ 必是 $-f(x)$ 的极小值点.　　(D) 对一切 $x$ 都有 $f(x)\leqslant f(x_0)$.

\vfill
\newpage
\qn{1991.9}
曲线 $y=\dfrac{1+\mathrm{e}^{-x^2}}{1-\mathrm{e}^{-x^2}}$（　　）
\\
\noindent (A) 没有渐近线.　　(B) 仅有水平渐近线.
\\
\noindent (C) 仅有铅直渐近线.　　(D) 既有水平渐近线又有铅直渐近线.

\vfill
\newpage
\qn{1991.10}
如图，$x$ 轴上有一线密度为常数 $\mu$，长度为 $l$ 的细杆，若质量为 $m$ 的质点到杆右端的距离为 $a$，已知引力系数为 $k$，则质点和细杆之间引力的大小为（　　）
\\
\noindent (A) $\int_{-l}^0\dfrac{km\mu}{(a-x)^2}\,\mathrm{d}x$.　　(B) $\int_0^l\dfrac{km\mu}{(a-x)^2}\,\mathrm{d}x$.
\\
\noindent (C) $2\int_{-\frac{l}{2}}^0\dfrac{km\mu}{(a+x)^2}\,\mathrm{d}x$.　　(D) $2\int_0^{\frac{l}{2}}\dfrac{km\mu}{(a+x)^2}\,\mathrm{d}x$.

\vfill
\newpage
\qn{1991.11}
设 $\begin{cases} x=t\cos t, \\ y=t\sin t, \end{cases}$ 求 $\dfrac{\mathrm{d}^2y}{\mathrm{d}x^2}$.

\vfill
\newpage
\qn{1991.12}
计算 $\int_1^4\dfrac{\mathrm{d}x}{x(1+\sqrt{x})}$.

\vfill
\newpage
\qn{1991.13}
求 $\lim\limits_{x\to0}\dfrac{x-\sin x}{x^2(\mathrm{e}^x-1)}$.

\vfill
\newpage
\qn{1991.14}
求 $\int x\sin^2 x\,\mathrm{d}x$.

\vfill
\newpage
\qn{1991.15}
求微分方程 $xy'+y=x\mathrm{e}^x$ 满足 $y(1)=1$ 的特解.
利用导数证明：当 $x>1$ 时，$\dfrac{\ln(1+x)}{\ln x}>\dfrac{x}{1+x}$.
求微分方程 $y''+y=x+\cos x$ 的通解.
曲线 $y=(x-1)(x-2)$ 和 $x$ 轴围成一平面图形，求此平面图形绕 $y$ 轴旋转一周所成的旋转体的体积.
如图，$A$ 和 $D$ 分别是曲线 $y=\mathrm{e}^x$ 和 $y=\mathrm{e}^{-2x}$ 上的点，$AB$ 和 $DC$ 均垂直 $x$ 轴，且 $|AB|:|DC|=2:1$，$|AB|<1$，求点 $B$ 和 $C$ 的横坐标，使梯形 $ABCD$ 的面积最大.
设函数 $f(x)$ 在 $(-\infty,+\infty)$ 上满足 $f(x)=f(x-\pi)+\sin x$，且 $f(x)=x$，$x\in[0,\pi)$. 计算 $\int_\pi^{3\pi}f(x)\,\mathrm{d}x$.

\vfill
\newpage
\end{document}